\section{Immutable Parent Pointers}
\label{sec:immutable-parent-pointers}

The integrity of the recursive spawning model rests upon an invariant that is both simple in statement and profound in consequence: the \emph{parent pointer} of any agent instance, once assigned at the moment of spawn, is immutable for the lifetime of that instance. This invariant pervades the entire architecture, cascading through authorization logic, chain traversal, termination conditions, and evidentiary reconstruction. Violating it would collapse the causal ancestry structure into an undirected graph, obliterating the directional traceability that is the architectural backbone of recursive spawning.

This section establishes the formal semantics of parent pointer immutability (\S\ref{sec:parent-pointer-semantics}), derives the chain formation recurrence and its base case (\S\ref{sec:chain-formation}), specifies the Fact Layer enforcement mechanism (\S\ref{sec:fact-layer-enforcement}), details the Purpose Layer spawn authorization protocol (\S\ref{sec:purpose-layer-authorization}), and presents the chain traversal algorithm with correctness guarantees (\S\ref{sec:chain-traversal-algorithm}).

% --- 4.1 Parent Pointer Semantics ---
\subsection{Parent Pointer Semantics}
\label{sec:parent-pointer-semantics}

\begin{definition}[Parent Pointer]
\label{def:parent-pointer}
Given two agent instances $p$ (the \emph{parent}) and $c$ (the \emph{child}), a \emph{ParentPointer} relation is established at spawn authorization and is formally defined as the binary predicate:

\[
\text{ParentPointer}(p, c) \coloneqq \texttt{TRUE}
\]

This relation is \emph{singular}: for any child $c$, there exists \emph{at most one} parent $p$ such that $\text{ParentPointer}(p, c) = \texttt{TRUE}$. The system enforces the constraint that $\nexists\; p_1 \neq p_2 : \text{ParentPointer}(p_1, c) \land \text{ParentPointer}(p_2, c)$. That is, no agent instance may have more than one parent.
\end{definition}

\begin{axiom}[Immutability Axiom]
\label{ax:parent-pointer-immutability}
Once $\text{ParentPointer}(p, c)$ is set to $\texttt{TRUE}$ at the moment of spawn authorization, it shall not be modified, cleared, redirected, or invalidated for the entire operational lifetime of $c$, regardless of subsequent system state, administrative action, or runtime condition.

\[
\forall\; c,\; \forall\; t > t_{\text{spawn}}(c),\; \text{ParentPointer}(p, c) = \texttt{TRUE} \implies \text{immutable}(p, c, t)
\]

where $t_{\text{spawn}}(c)$ denotes the timestamp of $c$'s spawn event.
\end{axiom}

The immutability axiom applies uniformly across all agent instance types in the recursive spawning hierarchy, including human agents, software agents, and hybrid human-agent composites. There is no override mechanism, no privilege escalation pathway, and no emergency bypass within the system's authorized operational envelope that permits retroactive modification of $\text{ParentPointer}$.

\begin{figure}[tb]
\centering
\fbox{
\begin{minipage}{0.92\textwidth}
\[
\text{Chain}(a) \;=\;
\begin{cases}
\varnothing & \text{if } \text{ParentPointer}(p, a) = \texttt{NULL} \quad \text{(root case)} \\[6pt]
\{\text{ParentPointer}(\text{parent}(a), a)\} \;\cup\; \text{Chain}(\text{parent}(a)) & \text{otherwise}
\end{cases}
\]
\vspace{4pt}
\par\noindent\textbf{Where:}
\begin{itemize itemtype="itemize"}
\item $\text{parent}(a)$ denotes the unique $p$ such that $\text{ParentPointer}(p, a) = \texttt{TRUE}$
\item $\text{ParentPointer}(\text{parent}(a), a)$ is a singleton representing the immediate ancestry link
\item $\text{Chain}(\text{parent}(a))$ is the recursive chain of the parent
\end{itemize}
\end{minipage}
}
\caption{Recursive definition of the ancestry chain for agent instance $a$}
\label{fig:chain-definition}
\end{figure}

The definition in Figure~\ref{fig:chain-definition} induces a well-founded recursion because the parent relationship is acyclic by construction: no agent can be its own ancestor (enforced by the spawn authorization protocol in \S\ref{sec:purpose-layer-authorization}). Consequently, the recursion terminates at the root human, which carries a null parent reference by definition.

% --- 4.2 Chain Formation ---
\subsection{Chain Formation: Recurrence and Base Case}
\label{sec:chain-formation}

The ancestry chain of an agent $a$ is defined recursively as shown in Figure~\ref{fig:chain-definition}. The construction is set-based and accumulates ancestry links through repeated application of the ParentPointer relation.

\begin{definition}[Ancestry Chain]
\label{def:ancestry-chain}
The \emph{Ancestry Chain} of an agent $a$, denoted $\text{Chain}(a)$, is the set of all ancestor links obtained by recursively applying the ParentPointer relation from $a$ to the root:

\[
\text{Chain}(a) \;=\; \bigcup_{i=0}^{n} \; \Bigl\{ \text{ParentPointer}(\text{ancestor}_i(a), \text{descendant}_i(a)) \Bigr\}
\]

where $\text{ancestor}_0(a) = a$, $\text{descendant}_0(a) = a$, and for $i > 0$, $\text{descendant}_i(a)$ is the parent of $\text{descendant}_{i-1}(a)$, until $\text{ParentPointer}(\text{parent}(a), a) = \texttt{NULL}$.
\end{definition}

The depth $n$ of the chain is the number of non-null ParentPointer links from $a$ to the root. Each link in the chain is a two-tuple $(p, c)$ representing the immutably recorded parent-child relationship.

\begin{property}[Chain Termination at Root Human]
\label{prop:chain-termination}
The root human agent, denoted $r$, is defined by the condition:

\[
\text{ParentPointer}(p, r) = \texttt{NULL} \quad \forall\; p
\]

For any agent $a$, iterating the chain formation recursively will eventually encounter a root human $r$ (or an agent with a null parent pointer), at which point the recursion terminates with the empty set as the base case.
\end{property}

\begin{note}
The root human represents the terminal node of the ancestry chain in the human-defined hierarchy. While it is theoretically possible for a purely synthetic agent lineage to have a synthetic root (e.g., a system-initialized agent with no parent), the recursive spawning model as defined herein assumes human initiation as the ultimate source. Any deviation from this assumption requires explicit redefinition of the spawn authorization semantics.
\end{note}

% --- 4.3 Fact Layer Enforcement ---
\subsection{Immutability Enforcement: The Fact Layer}
\label{sec:fact-layer-enforcement}

The Fact Layer is the substrate upon which the immutability of parent pointers is operationally enforced. It serves as the authoritative record-keeper for all spawned agents and their parent-child relationships, and it is architecturally isolated from the authorization and purpose-resolution logic to prevent cross-layer interference with ancestry records.

\begin{architecture}[Fact Layer]
\label{arch:fact-layer}
The Fact Layer consists of:
\begin{enumerate}
\item \textbf{Ancestry Registry}: A persistent, append-only data structure that records every spawn event as an immutable tuple $(\text{timestamp}, \text{parent\_id}, \text{child\_id}, \text{authorization\_proof}, \text{spawn\_context})$. Once written, no tuple is ever modified or deleted.

\item \textbf{Parent Pointer Store}: A key-value store mapping each $\text{child\_id}$ to its unique $\text{parent\_id}$, initialized at spawn time and thereafter read-only. Lookups are $O(1)$ on average.

\item \textbf{Write-Once Enforcement}: The Fact Layer enforces a write-once policy at the architectural level. Any code path that attempts to modify an existing parent pointer entry is rejected at the syscall boundary. The enforcement is implemented via immutable reference semantics in the underlying storage layer and is not dependent on runtime permission checks.
\end{enumerate}
\end{architecture}

\begin{invariant}[Fact Layer Integrity]
\label{inv:fact-layer-integrity}
For any agent $c$ that has been spawned within the system:

\[
\text{parent}(c) = p \iff \text{FactLayer.RegistryContains}(\text{tuple}(p, c)) \land \text{immutable}(\text{tuple}(p, c))
\]

The Fact Layer never returns a null parent for a non-root agent unless the agent was explicitly spawned with a null parent designation (reserved for root humans and system-initialized seed agents).
\end{invariant}

The Fact Layer exposes a narrow, well-defined interface to other layers:

\begin{itemize}
\item \texttt{FactLayer.GetParent(child\_id)} $\rightarrow$ parent\_id \quad [read-only query]
\item \texttt{FactLayer.RecordSpawn(parent\_id, child\_id, context)} $\rightarrow$ \texttt{BOOL} \quad [write-once at spawn only]
\item \texttt{FactLayer.GetChainDepth(agent\_id)} $\rightarrow$ depth \quad [derived read-only query]
\item \texttt{FactLayer.ValidateChain(agent\_id)} $\rightarrow$ \texttt{BOOL} \quad [acyclicity verification]
\end{itemize}

No update, delete, or redirect operations are exposed or permitted. The Fact Layer is a \emph{read-mostly} data structure with a single controlled write path that is invoked exclusively during the spawn authorization phase.

% --- 4.4 Purpose Layer Authorization ---
\subsection{Spawn Authorization: The Purpose Layer}
\label{sec:purpose-layer-authorization}

The Purpose Layer is the sole authority for permitting a spawn event. It resolves the question of whether a given parent agent is authorized to spawn a child agent, and it does so by evaluating the declared purpose of the spawn against the agent's operational mandate, competency scope, and authorization level.

\begin{protocol}[Spawn Authorization Protocol]
\label{proto:spawn-authorization}
\begin{enumerate}
\item The parent agent $p$ submits a spawn request to the Purpose Layer, declaring:
    \begin{itemize}
    \item $\text{Intent}(p, c)$: the intended purpose of spawning $c$
    \item $\text{Scope}(c)$: the operational scope of the proposed child
    \item $\text{AuthorizationLevel}(p)$: the parent agent's current authorization tier
    \end{itemize}

\item The Purpose Layer evaluates whether $\text{Intent}(p, c)$ falls within the authorized scope of $p$.

\item If the evaluation succeeds, the Purpose Layer issues an \texttt{AuthorizationToken}$(p, c, \text{timestamp})$, which is a cryptographically signed assertion that the spawn is purpose-compliant.

\item The Fact Layer receives the AuthorizationToken and records the spawn event, irrevocably establishing $\text{ParentPointer}(p, c) = \texttt{TRUE}$.

\item If the evaluation fails, no spawn record is created, and $\text{ParentPointer}$ remains undefined for $c$.
\end{enumerate}
\end{protocol}

The Purpose Layer enforces two critical constraints that indirectly strengthen the immutability of parent pointers:

\begin{constraint}[No Retroactive Parent Assignment]
A spawn request for agent $c$ must be submitted and authorized \emph{before} $c$ is instantiated. The Fact Layer will not accept a parent pointer assignment for an agent that already exists without a corresponding authorized spawn event. This prevents the scenario where a parent pointer is assigned to an already-active agent after the fact, which would bypass the Purpose Layer's authorization evaluation.

\item \textbf{Constraint 2 — Unidirectional Authorization Flow}
Spawn authorization is strictly top-down: a parent authorizes a child. There is no mechanism for a child to request or compel its own parent assignment post-spawn. This prevents self-referential parent assignments and ensures the ancestry structure remains a directed acyclic graph (DAG).
\end{constraint}

The Purpose Layer also enforces the acyclicity constraint indirectly: since each spawn request references the parent as an already-authorized agent, and the parent is already recorded in the Fact Layer, the system can verify that the proposed child is not an ancestor of the parent before issuing authorization. This prevents circular spawn chains.

% --- 4.5 Chain Traversal Algorithm ---
\subsection{Chain Traversal Algorithm}
\label{sec:chain-traversal-algorithm}

Given the immutable parent pointer structure and the recursively defined chain, we can derive an algorithm for traversing the ancestry chain of any agent $a$, terminating at the root human (null parent).

\begin{algorithm}[Ancestry Chain Traversal]
\label{alg:chain-traversal}
\begin{algorithmic}
\REQUIRE Agent instance identifier $a$
\ENSURE Ordered list of ancestry links from $a$ to root (inclusive)
\STATE $\text{chain} \leftarrow \text{empty list}$
\STATE $\text{current} \leftarrow a$
\WHILE{$\text{current} \neq \texttt{NULL}$}
    \STATE $\text{parent} \leftarrow \text{FactLayer.GetParent}(\text{current})$
    \IF{$\text{parent} = \texttt{NULL}$}
        \STATE $\text{APPEND } (\text{NULL}, \text{current})$ \COMMENT{Root link: null parent of root agent}
        \STATE $\text{BREAK}$
    \ENDIF
    \STATE $\text{APPEND } (\text{parent}, \text{current})$
    \STATE $\text{current} \leftarrow \text{parent}$
\ENDWHILE
\RETURN $\text{chain}$
\end{algorithmic}
\end{algorithm}

The correctness of Algorithm~\ref{alg:chain-traversal} follows from the immutability and acyclicity invariants. Since $\text{FactLayer.GetParent}(x)$ always returns the unique, immutable parent of $x$ (or NULL if $x$ is a root), and since the parent relation is acyclic, the loop is guaranteed to terminate. The termination is guaranteed by the well-foundedness of the parent relation: there is no infinite descending chain because the maximum depth is bounded by the number of agents in the system, and acyclicity guarantees no loop.

\begin{property}[Algorithm Correctness]
\label{prop:traversal-correctness}
Algorithm~\ref{alg:chain-traversal} satisfies the following:
\begin{enumerate}
\item \textbf{Termination}: The loop executes at most $n+1$ iterations where $n$ is the chain depth of $a$, and $n$ is finite due to acyclicity.
\item \textbf{Completeness}: Every ancestor of $a$ appears exactly once in $\text{chain}$.
\item \textbf{Order Correctness}: The chain is ordered from $a$ (deepest descendant) to the root (null parent terminal node).
\item \textbf{Idempotence}: Running the algorithm multiple times on the same $a$ yields identical output, due to parent pointer immutability.
\end{enumerate}
\end{property}

The traversal algorithm supports several critical operations across the recursive spawning system:

\begin{itemize}
\item \textbf{Evidentiary Reconstruction}: Given a child agent $c$, the full ancestry chain can be reconstructed to establish accountability and provenance.
\item \textbf{Authorization Propagation Analysis}: The chain can be traversed to determine which ancestor authorized which descendant link, supporting audit and compliance verification.
\item \textbf{Competency Scope Inheritance Verification}: Each hop in the chain represents a point where a Purpose Layer authorization was granted; traversing the chain validates that each spawn in the lineage was properly authorized.
\item \textbf{Root Identification}: The traversal algorithm terminates at the root human, which is the ultimate source of authorization for the entire chain.
\end{itemize}

% --- 4.6 Implications ---
\subsection{Implications and Architectural Consequences}
\label{sec:implications}

The immutability of parent pointers is not merely a data integrity constraint — it is the foundation upon which the entire recursive spawning model operates. Several architectural consequences follow directly from this invariant:

\begin{enumerate}
\item \textbf{Auditability}: Every agent has a complete, tamper-evident ancestry chain. Since parent pointers cannot be modified after spawn, the chain is a reliable record of causal relationships. This has direct implications for compliance, accountability, and dispute resolution.

\item \textbf{Deterministic Provenance}: Given any agent $c$, the complete provenance chain $\text{Chain}(c)$ can be computed at any time and will always produce the same result. This reproducibility is essential for AI oversight systems that require reliable historical reconstruction.

\item \textbf{Causal Clarity}: The immutable parent pointer eliminates ambiguity about the origin of any agent. There is exactly one path from any agent to the root, and that path cannot be altered after the fact.

\item \textbf{Formal Verification}: The immutability invariant admits formal verification within a theorem-proving framework. Properties such as acyclicity, termination, and authorization completeness can be proven given the immutability axiom as a starting assumption.

\item \textbf{System Robustness}: Since parent pointers cannot be corrupted or redirected by adversarial action (within the authorized operational envelope), the system is resilient against certain classes of attacks that attempt to falsify ancestry records.
\end{enumerate}

The recursive spawning model therefore depends on the ironclad enforcement of parent pointer immutability as a non-negotiable architectural invariant. Any relaxation of this constraint — whether through design oversight, implementation error, or adversarial compromise — would fundamentally undermine the integrity of the ancestry chain and the trustworthiness of the entire agent ecosystem.

\end{document}